\documentclass[11pt]{article}
\usepackage[margin=1in]{geometry}
\usepackage{amsmath,amssymb}
\usepackage{booktabs}
\usepackage{longtable}
\usepackage{graphicx}
\usepackage{hyperref}
\usepackage{xcolor}
\usepackage[numbers,sort&compress]{natbib}
\usepackage{microtype}
\graphicspath{{../figures/}}
\newcommand{\todo}[1]{\textcolor{red}{[TODO: #1]}}
\newcommand{\bigO}{\mathcal{O}}

\title{When Does a Theoretically Faster Algorithm Actually Win?\\
\large An empirical crossover atlas for four classic algorithm pairs}
\author{Dhanush Paturu \\ Nitte Meenakshi Institute of Technology, Bengaluru, India \\ \href{https://orcid.org/0009-0004-4918-6181}{ORCID 0009-0004-4918-6181}}
\date{\today}

\begin{document}
\maketitle

\begin{abstract}
Asymptotic analysis promises that an algorithm with a better growth rate will
eventually outperform a simpler one, but says nothing about \emph{when}. For four
classic pairs---radix sort versus introsort, Strassen versus naive matrix
multiplication, Karatsuba versus schoolbook integer multiplication, and
Dijkstra's algorithm with a Fibonacci heap versus a binary heap---we measure the
input size at which the asymptotically faster algorithm first wins in wall-clock
time on a current commodity CPU, and how that crossover moves with input
structure (key width, sortedness, graph density) and with the tuning parameter
every recursive algorithm has (its base-case cutoff). Both sides of each pair are
plain, single-threaded C++ compiled with the same flags, and the recursive
algorithms bottom out in exactly the kernel they are compared against.
Crossovers are estimated on geometric size grids with bootstrap confidence
intervals over interleaved repetitions. On an Intel Core i7-14650HX, radix sort overtakes introsort at about 1\,500
keys regardless of key width, but never on nearly-sorted or reverse-sorted
input, where introsort stays $3.7\times$ and $8\times$ ahead at $2^{26}$ keys; Strassen
overtakes the naive triple loop from $n=32$ on dimensions rich in factors of
two but only from $n\approx300$ on arbitrary dimensions, and the loop order of
the \emph{naive} baseline is worth more than Strassen's exponent; Karatsuba
overtakes schoolbook multiplication at 34 limbs, within a factor of $1.3$ of
the threshold a tuned library uses on the same core family; and the Fibonacci heap never beats a binary heap on any sparse graph,
reaching parity only on complete graphs at the edge of memory while still
losing to the lazy-deletion binary heap that practitioners actually write.
Crossovers are reported with bootstrap intervals and classified as sharp,
gradual or marginal, and we argue that a wide interval is itself the practical
answer: a range over which the choice does not matter. The result is a practical reference: a map of where the constants
hidden by $\bigO$-notation stop mattering, with the code and data to redraw it
on any machine.
\end{abstract}

\section{Introduction}

The standard argument for preferring an algorithm with a better asymptotic
bound is that, for large enough inputs, it must win. The argument is correct and
incomplete: it is silent about the size at which ``large enough'' begins, and in
practice that size is set by constant factors that the analysis discards---the
number of passes over memory, the size of the working set relative to the cache
hierarchy, the predictability of branches, and the overhead of recursion and
bookkeeping. Experienced programmers carry folk knowledge about these
crossovers (``radix sort is only worth it for big arrays'', ``nobody uses
Fibonacci heaps''), but the folk knowledge is rarely written down with numbers,
and when it is, the numbers are a decade old and from a different cache
hierarchy.

This paper measures the crossovers directly. For each of four algorithm pairs
we sweep the input size over two and a half to seven orders of magnitude, vary the
structure of the input and the tuning parameter of the faster algorithm, and
record the first size at which the median wall-clock time of the theoretically
faster algorithm drops below that of the simpler one. We call the resulting
table and plots a \emph{crossover atlas}.

\paragraph{Contributions.}
\begin{enumerate}
  \item A controlled, reproducible measurement of crossover points for radix
  sort vs.\ introsort, Strassen vs.\ naive matrix multiplication, Karatsuba vs.\
  schoolbook multiplication and Fibonacci-heap vs.\ binary-heap Dijkstra, with
  bootstrap confidence intervals (Section~\ref{sec:results}).
  \item An analysis of \emph{what moves} each crossover: key width and
  sortedness for sorting, the base-case cutoff for the two recursive
  algorithms, and graph density and the lazy-deletion idiom for shortest
  paths (Section~\ref{sec:results}).
  \item An account of the mechanism behind each crossover---passes over
  memory, the access pattern of the base case, pointer chasing against
  arrays---argued from working-set sizes and logged operation counts
  (Section~\ref{sec:mechanisms}).
  \item A benchmarking harness that handles the practical difficulties
  (sub-microsecond inputs, cubic baselines that would otherwise dominate the
  sweep, drift between repetitions) and that re-draws the atlas on any Linux
  machine with one command (Section~\ref{sec:method}, repository in
  Section~\ref{sec:repro}).
\end{enumerate}

\section{Background and related work}\label{sec:background}

\subsection{Experimental algorithmics}
The case for measuring algorithms rather than only analysing them was made
systematically by Johnson~\cite{johnson2002theoretician}, whose guide for
theoreticians lists the pitfalls this paper tries to avoid (uninteresting
instances, unfair baselines, unreported machines), and by
McGeoch~\cite{mcgeoch2012guide}, whose book is the standard reference on
experimental design for algorithms. Sanders~\cite{sanders2009algorithm} frames
the activity as \emph{algorithm engineering}: a cycle of design, analysis,
implementation and experiment in which the constant factors are first-class
objects. The measurement side has its own literature of cautionary results.
Mytkowicz et al.~\cite{mytkowicz2009producing} showed that link order and
environment-variable size alone can change measured runtimes by tens of
percent, more than many published effects; Curtsinger and
Berger~\cite{curtsinger2013stabilizer} proposed randomising memory layout so
that such effects average out; and Hoefler and Belli~\cite{hoefler2015scientific}
surveyed reporting practice and argued for repetition, interval estimates and
full disclosure of the machine. LaMarca and Ladner~\cite{lamarca1999influence}
is the direct ancestor of this study for sorting: by instrumenting cache
misses they showed that the instruction count of a sorting algorithm is a poor
predictor of its time once the input leaves cache, and that memory-aware
variants of heapsort and mergesort outrun the textbook versions.

\subsection{The four pairs}
\paragraph{Sorting.}
Introsort~\cite{musser1997introspective}---quicksort with a heapsort fallback
and insertion sort for small ranges---is what \texttt{std::sort} runs in every
major C++ standard library, which makes it the baseline a practitioner actually
faces. Rahman and Raman~\cite{rahman2001adapting} adapted LSD radix sort to the
cache and TLB and showed it beating comparison sorts on large arrays of
integers; Satish et al.~\cite{satish2010fast} compared radix and merge sort
across CPUs and GPUs and found radix sort faster for 32-bit keys but merge sort
preferable for wide keys and on bandwidth-limited hardware; Polychroniou and
Ross~\cite{polychroniou2014comprehensive} studied the partitioning step that
both families share at billion-key scale. None of these report the crossover
size against \texttt{std::sort} as an estimated quantity, and all use SIMD
and multithreading that move it.

\paragraph{Matrix multiplication.}
Strassen's 1969 algorithm~\cite{strassen1969gaussian} is the archetype of an
asymptotic improvement that is practical only after engineering.
Huss-Lederman et al.~\cite{husslederman1996implementation} gave the first
careful implementation study, establishing that the recursion must stop at a
tuned cutoff and that odd dimensions need peeling or padding; D'Alberto and
Nicolau~\cite{dalberto2007adaptive} tuned that cutoff adaptively to the memory
hierarchy; Benson and Ballard~\cite{benson2015framework} showed that the best
fast algorithm depends on matrix shape; and Huang et
al.~\cite{huang2016strassen}, by fusing Strassen's additions into a BLIS
kernel, brought the \emph{switch point}---the dimension above which one level of
Strassen pays off against a vendor-grade DGEMM---down from the roughly 2\,000
of earlier implementations to as little as 500 for their best variant, on a
single Ivy Bridge Xeon core. Those crossovers are against the strongest
possible baseline. Ours is against the baseline a
programmer writes without a BLAS, and is correspondingly much earlier.

\paragraph{Integer multiplication.}
Karatsuba and Ofman's algorithm~\cite{karatsuba1962multiplication} is the
first sub-quadratic multiplication and the first step of the tower that ends,
for now, at Harvey and van der Hoeven's $\bigO(n\log n)$
algorithm~\cite{harvey2021integer}, which is correct and
unimplementable. Bernstein~\cite{bernstein2001multidigit} gives the modern
unified view. The GNU Multiple Precision library~\cite{gmp2023manual} switches
from schoolbook to Karatsuba at a machine-tuned threshold
(\texttt{MUL\_TOOM22\_THRESHOLD}), which its manual notes can be as low as ten
64-bit limbs; the parameter file GMP ships for Alder Lake sets it to 13
limbs~\cite{gmp2022alderlake}. This gives an external reference point for our
measured crossover.

\paragraph{Shortest paths.}
Fredman and Tarjan~\cite{fredman1987fibonacci} introduced the Fibonacci heap
precisely to improve Dijkstra's algorithm from $\bigO(m\log n)$ to
$\bigO(m+n\log n)$, and the data structure has been a textbook fixture ever
since; Brodal, Lagogiannis and Tarjan~\cite{brodal2012strict} later made the
bounds worst-case. The empirical record is uniformly unkind to it. Cherkassky,
Goldberg and Radzik~\cite{cherkassky1996shortest} found simple heaps and bucket
structures faster on every graph family they tried; Chen et
al.~\cite{chen2007priority} dropped the Fibonacci heap from their comparison
after preliminary runs and found that Dijkstra variants \emph{without}
decrease-key were at least $1.4\times$ faster than those with it on random
graphs of average degree 8 (and faster on every road network they tried); and
Larkin, Sen and Tarjan~\cite{larkin2014backtobasics}, on the full USA road
network, measured the Fibonacci heap at $3.15\times$ the time of an implicit
4-ary heap, with one decrease-key for every fourteen inserts, and found the L1
cache-miss rate to be the counter most strongly correlated with wall-clock
time. Our results
agree with all of these and add the one regime, complete graphs, in which the
asymptotic argument finally holds.

\subsection{Hidden constants and galactic algorithms}
Lipton~\cite{lipton2010galactic} coined \emph{galactic algorithm} for
one whose asymptotic advantage appears only on inputs larger than any that
will be run, with fast matrix multiplication as the canonical example: the
exponent has been pushed to $\omega<2.372$~\cite{alman2024refined,
alman2024asymmetry}, yet nothing past Strassen is used in practice. This paper
sits at the opposite end of that spectrum. Its four pairs are algorithms that
are actually used, and the question is not whether their crossovers exist but
where, with what precision, and what moves them. The hardware background that
answers ``why''---cache levels, memory-level parallelism, branch
prediction---is laid out by Drepper~\cite{drepper2007memory} and Fog~\cite{fog2026optimizing}; the
Linux \texttt{perf\_event\_open} interface~\cite{linux2026perfeventopen} is
how the harness reads the counters where it can.


\section{Method}\label{sec:method}

\subsection{Algorithm pairs and what is held fixed}

Each pair consists of a \emph{faster-in-theory} algorithm $F$ and a
\emph{simpler} algorithm $S$ (Table~\ref{tab:pairs}). Three rules keep the
comparison about the algorithmic idea rather than about implementation effort.
\begin{enumerate}
  \item \textbf{Same language, compiler and flags.} C++17, GCC, \texttt{-O2},
  no \texttt{-march=native} (the x86-64 baseline ISA), no hand-written SIMD, no
  threads. (The harness can rebuild with \texttt{-O3 -march=native} for a
  secondary experiment; that experiment is not part of this paper.)
  \item \textbf{Same base case.} The recursive algorithms (Strassen, Karatsuba)
  bottom out in exactly the kernel they are compared against, so the measured
  difference is attributable to the recursive structure, its extra additions
  and its memory traffic.
  \item \textbf{Competent, textbook implementations on both sides.} Radix sort
  computes all digit histograms in one pass and skips passes in which every key
  shares a digit; Karatsuba uses a caller-provided scratch buffer; the
  Fibonacci heap pre-allocates one node per vertex. The simpler algorithms are
  what a careful programmer writes: \texttt{std::sort}, the $i$-$k$-$j$ loop
  order, 64-bit multiply-accumulate, an indexed binary heap with
  decrease-key---and, as a second baseline, the lazy \texttt{std::priority\_queue}
  idiom that is far more common in practice.
\end{enumerate}

\begin{table}[t]
\centering\small
\caption{The four algorithm pairs and the variables swept.}
\label{tab:pairs}
\begin{tabular}{p{2.2cm}p{3.3cm}p{3.3cm}p{5.6cm}}
\toprule
Pair & $F$ & $S$ & Swept \\
\midrule
Sorting & LSD radix sort, $\bigO(n\cdot 64/d)$ for $d$-bit digits & introsort (\texttt{std::sort}), $\bigO(n\log n)$ & $n\in[16,2^{26}]$; key width 8/16/32/64 bits; nearly-sorted, reverse, few-unique, organ-pipe; $d\in\{8,11,16\}$ \\
Matrix multiply & Strassen, $\bigO(n^{2.807})$ & naive triple loop, $\bigO(n^3)$ & $n\in[16,4096]$ incl.\ non-powers of two; cutoff $\in\{8,\dots,256\}$; $i$-$k$-$j$ vs $i$-$j$-$k$ baseline \\
Integer multiply & Karatsuba, $\bigO(n^{1.585})$ & schoolbook, $\bigO(n^2)$ & $n\in[8,2^{20}]$ 32-bit limbs; cutoff $\in\{8,\dots,128\}$; random and all-ones limbs \\
Shortest paths & Dijkstra + Fibonacci heap, $\bigO(m+n\log n)$ & Dijkstra + binary heap, $\bigO((n+m)\log n)$ & $n\in[10^2,8.4\times10^6]$; average degree $\in\{2,\dots,256\}$; random / 2-D grid / complete graphs; lazy baseline \\
\bottomrule
\end{tabular}
\end{table}

\subsection{Inputs}
Inputs are generated deterministically from a 64-bit seed combined with $n$
(xoshiro256** seeded through splitmix64), so every cell of the experiment is
reproducible from the command line that produced it. Sizes form a geometric
grid (factor 2 for the main sweeps, 1.5 for integer multiplication, and a
hand-picked list for matrices that includes sizes such as 24, 48 and 96 that
exercise Strassen's padding path). A \emph{cell} is a (pair, input family,
parameter, $n$) tuple. Cells that make no sense are skipped by rule (a complete
graph on $10^7$ vertices); an algorithm whose single run exceeds a wall-clock
budget is dropped for larger $n$ in that series, so that cubic baselines do not
dominate the running time of the experiment.

\subsection{Measurement}
\paragraph{Timer.} \texttt{std::chrono::steady\_clock} (\texttt{CLOCK\_MONOTONIC}
through the vDSO, about 20\,ns overhead). For inputs whose single run exceeds
2\,ms the input is copied outside the timed region and only the algorithm is
timed. Shorter runs are batched: the harness runs $K$ iterations of (copy
input; run algorithm) and reports the per-iteration mean, with $K$ chosen from
a warm calibration run; because in-place algorithms need a fresh copy per
iteration, the copy is inside the timed region, and its cost is measured
separately with the same $K$ on the same pre-allocated buffer and subtracted
in the analysis (never below 5\% of the raw time).

\paragraph{Repetitions and interleaving.} Each cell is measured $r=5$ times
after $w=2$ untimed warm-up runs. Within the repetition loop the algorithms
alternate in a rotating order, so slow drift in machine state (thermal
throttling, background activity, frequency changes) is shared equally by $F$
and $S$. The per-cell statistic is the median; the interquartile range is
retained and plotted.

\paragraph{Hardware counters.} When \texttt{perf\_event\_open} is permitted,
each measurement also records cycles, retired instructions, L1d read misses,
last-level-cache misses and branch mispredictions for the user-space portion of
the run, scaled for multiplexing. These are explanatory, not primary: the
crossover is defined on wall-clock time. Where counters are unavailable (virtual
machines, WSL2) the study is timing-only and says so.

\paragraph{Correctness.} Every benchmark binary cross-checks the outputs of $F$
and $S$ on the last repetition of each cell and aborts on mismatch. Separate
randomized test suites compare each implementation with an independent
reference (\texttt{std::sort}; the $i$-$j$-$k$ product; schoolbook plus the
modular identity $(ab)\bmod p = ((a\bmod p)(b\bmod p))\bmod p$ with
$p=2^{61}-1$; Bellman--Ford), and both heaps are exercised directly against a
scan-for-minimum oracle under random push, decrease-key and pop sequences.

\subsection{Runs and a measurement correction}\label{sec:method-patch}
The reported dataset is the union of three runs made in one session on the
same machine, with the same binary flags: \texttt{full} (the sweeps of
Table~\ref{tab:pairs}, 5.3 hours), \texttt{refine} (factor-1.19 grids over
the neighbourhood of each crossing, 20 minutes) and \texttt{patch} (70
minutes). The patch exists because the first version of the harness measured
the input-copy baseline into a freshly allocated buffer, so that for large
inputs the baseline included the operating system's page-fault cost while the
copy inside the timed loop, reusing the buffer, did not; the subtraction then
removed too much. The effect is below 5\% wherever the copy is small relative
to the run, which is almost everywhere, but for complete graphs (470\,MB of
adjacency at 7\,700 vertices against a 100\,ms run) and for the largest
arrays sorted by the fastest configurations it was large enough to invalidate
the cell. We fixed the harness as described above, re-measured every series in
which the copy had exceeded 5\% of the raw time (dense graphs at all sizes,
arrays of $2^{21}$ keys and above, sparse graphs of $2^{20}$ vertices and
above, and a handful of small cells), and let the re-measurement
\emph{replace} the original cells over the range of $n$ it covers rather than
pool with them. Both the original and the corrected data are in the
repository; the analysis script's \texttt{--supersede} flag applies the
replacement.

\subsection{Machine}
All reported measurements were made on a laptop with an Intel Core
i7-14650HX (Raptor Lake, 8 performance cores + 8 efficiency cores, 24
threads; 48\,KB L1d and 2\,MB L2 per performance core, 30\,MB shared L3)
running Ubuntu under WSL2 (Linux 6.18, Microsoft hypervisor) with 7.6\,GB of
memory visible to the guest, GCC 15.2.0 with \texttt{-O2}, CMake 4.2. The
laptop was on mains power in the manufacturer's \emph{Performance} mode, the
Windows power plan was set to \emph{High performance} with maximum processor
state 99\% (which disables Turbo Boost on Intel parts), the WSL virtual
machine was restricted to the performance cores through its process
affinity, and the benchmark thread was pinned to one virtual CPU. The hypervisor hides
hardware performance counters, so the reported run is timing-only; the
harness collects counters automatically on bare-metal Linux, where it also
sets the performance governor, disables turbo and pins the frequency
(\texttt{scripts/setup\_machine.sh}). We report the machine exactly and treat
absolute times as specific to it; the crossover \emph{ordering} and
\emph{order of magnitude} are the transferable results.

\subsection{Crossover estimation}
For a series (pair, family, parameter) let $\rho(n)=\tilde t_F(n)/\tilde t_S(n)$
be the ratio of median runtimes on the size grid. The crossover $n^*$ is the
\emph{sustained} crossing: the grid interval after which $\rho<1$ at every
larger measured size, interpolated linearly in $(\log n,\log\rho)$. A
transient dip below 1 followed by a return above it is not counted; it remains
visible in the ratio plots and in the reported minimum ratio. Series with $\rho<1$ already at the smallest $n$ are
reported as \emph{always faster} (from $n_{\min}$) and series with $\rho\ge1$
throughout as \emph{never faster} (up to $n_{\max}$), together with
$\rho(n_{\max})$ so the reader can see how far from crossing they are.
Uncertainty comes from a non-parametric bootstrap: the $r$ repetitions of every
cell are resampled with replacement (1000 replicates), medians and $n^*$ are
recomputed, and the shortest interval containing 95\% of the replicates (in
$\log n$) is reported. The shortest rather than the central interval is used
because the sustained-crossing statistic can be bimodal: a resample in which
one noisy mid-range cell pops above $\rho=1$ relocates the crossing to the far
end of the sweep, and a percentile interval would stretch over that tail even
when nearly all of the mass sits tightly around the true crossing.

Two caveats govern how the intervals should be read. First, $n^*$ is
interpolated between grid points, so an interval cannot resolve a crossing
more finely than the grid allows; the main sweeps use a factor-2 grid and the
\texttt{refine} profile re-sweeps the neighbourhood of each crossing on a
factor-1.19 grid for that purpose. Second, and more important, the width of an
interval is itself a result. We classify each crossing as \emph{sharp} when
its interval is narrower than one grid step or than a factor of 1.5,
whichever is wider (a $\pm20\%$ interval is sharp by any practical standard,
even on a fine grid), as \emph{gradual} otherwise, and as \emph{marginal}
when the faster algorithm is never more than 10\% ahead by the largest size
measured, whatever the interval. A gradual crossing means that $\rho$ approaches
1 slowly, i.e.\ the two algorithms are within measurement noise of each other
across the whole interval; the honest statement is then not ``the fast
algorithm wins from $n^*$'' but ``over this range it makes no measurable
difference, and beyond it the fast algorithm wins by the ratio we report at
$n_{\max}$''. Gradual crossings are reported as ranges, never as points. Empirical scaling exponents are fitted as $\log t = b\log n + c$
over the largest decade of $n$ for each algorithm, to show how far the
measured growth is from the asymptotic bound in the measured range.


\section{Results}\label{sec:results}

All numbers in this section come from the \texttt{full} profile (factor-2 size
grids, 5 timed repetitions after 2 warm-ups, 2\,411 cells, 5.3 hours) on the
machine of Section~\ref{sec:method}, supplemented by the \texttt{refine}
profile (factor-1.19 grids around the crossings) and the \texttt{patch}
profile described in Section~\ref{sec:method-patch}.
Table~\ref{tab:crossovers} gives every crossing with its interval and class.

\subsection{Sorting: radix vs.\ introsort}

\begin{figure}[t]
  \centering
  \includegraphics[width=\linewidth]{fig_runtime_sorting}
  \caption{Median runtime of LSD radix sort (11-bit digits) and introsort by
  input family. Dotted lines mark the $n$ at which the working set
  ($16n$ bytes) equals L1d, L2 and L3. Radix sort sits on a flat floor
  below $n\approx10^3$, where the histogram pass dominates.}
  \label{fig:runtime-sorting}
\end{figure}
\begin{figure}[t]
  \centering
  \includegraphics[width=0.72\linewidth]{fig_ratio_sorting_by_dist}
  \caption{Runtime ratio radix/introsort by input family (11-bit digits).
  Uniform keys of four widths cross between 1.4\,k and 1.7\,k; nearly-sorted
  and reverse-sorted input never cross.}
  \label{fig:ratio-sorting}
\end{figure}

On uniformly random keys radix sort wins from about 1\,500 elements and the
crossing is sharp (Figures~\ref{fig:runtime-sorting}
and~\ref{fig:ratio-sorting}): 1.43\,k keys for 8-bit values, 1.60\,k for
16-bit, 1.51\,k for 32-bit and 1.71\,k for 64-bit, each with an interval of
$\pm10\%$ or less. The near-independence from key width is itself a finding:
with 11-bit digits a 64-bit key needs six passes and an 8-bit key one, yet the
crossing barely moves, because below $\sim10^3$ keys both algorithms are
dominated by fixed costs---the histogram and prefix-sum setup on one side,
insertion-sort leaves on the other---and above it both scale as $n$ to within
a few percent over the whole range (fitted exponents 1.00--1.06 for radix,
1.02--1.14 for introsort). Above the crossing the advantage settles between
$2.7\times$ (64-bit keys: $2.44$\,s against $6.63$\,s at $2^{26}$) and
$4\times$ (16-bit keys: $1.14$\,s against $4.55$\,s).

Digit width is the practitioner's knob and it is not free. Eleven-bit digits
are best at every key width; eight-bit digits cost one or two extra passes
and move the 64-bit crossing only to 1.67\,k, but sixteen-bit digits are a
mistake for wide keys: the $2^{16}$-counter histogram per pass is 512\,KB,
larger than L1d, and at 64-bit keys the crossing moves to 7\,k and the
asymptotic advantage shrinks to $1.4\times$ (the gradual interval reaching
8.6\,M in Appendix~\ref{app:table} reflects the ratio hovering between 0.5
and 0.75 across the whole range).

The structured families overturn the verdict. With only 100 distinct values
introsort's partitioning handles duplicates well and the crossing moves to
2.7\,k with an eventual advantage of only $2\times$. On nearly-sorted input
radix sort \emph{never} wins: introsort is $3.7\times$ faster at $2^{26}$ keys
($1.15$\,s vs $4.26$\,s). On reverse-sorted input it is $8\times$ faster
($0.60$\,s vs $4.92$\,s)---median-of-three pivot selection turns a reversed
array into perfectly balanced partitions with perfectly predictable branches,
and introsort runs $11\times$ faster on it than on random data of the same
size. The organ-pipe input, ascending then descending, is the intermediate
case: radix sort crosses at 1.6\,k but only ever gets $1.4\times$ ahead.
Radix sort performs the same scattered writes whether the input is sorted or
not; an algorithm that is oblivious to structure cannot exploit it.

\subsection{Matrix multiplication: Strassen vs.\ naive}

\begin{figure}[t]
  \centering
  \includegraphics[width=0.85\linewidth]{fig_heatmap_matmul}
  \caption{$\log_2$ of the runtime ratio Strassen/naive ($i$-$k$-$j$) over
  matrix size and base-case cutoff on the power-of-two-rich grid; blue means
  Strassen wins.}
  \label{fig:heatmap-matmul}
\end{figure}

Strassen's algorithm is practical much earlier than folklore suggests, with
one important qualification about matrix shape. On sizes with many factors of
two ($n = 2^k$ and $3\cdot2^k$, the grid every benchmark uses) and a base-case
cutoff of 32, it beats the $i$-$k$-$j$ triple loop at every size from $n=32$
up---$0.83\times$ at $n=64$, $0.71\times$ at $256$, $0.39\times$ at $2048$,
$0.31\times$ at $4096$ ($13.4$\,s vs $43.3$\,s)---and the fitted exponent
over the top decade is $2.78$ against the theoretical $2.81$, with the naive
loop at $3.21$.

The refine grid, which includes sizes such as 129, 154, 217 and 259, tells a
different story. Our implementation handles an odd dimension at any level of
the recursion by padding that level to even size, which is the textbook
approach; an $n$ with few factors of two is padded at level after level, and
every padding copies both operands. On those sizes Strassen with cutoff 32 is
up to $1.2\times$ \emph{slower} than naive until $n\approx260$, and the
sustained crossing over all measured sizes is $n\approx300$ (interval
$[231, 308]$), five times later than on the composite grid. Both numbers are
right; which one applies depends on whether the caller controls the matrix
dimensions. Production implementations pad or peel once at the top level,
which would remove this effect, and the comparison shows how much it is worth.

The cutoff behaves as expected (Figure~\ref{fig:heatmap-matmul}): 8 is one
level too deep and costs a factor of $1.6$ at $n=4096$ ($21.2$\,s vs
$13.4$\,s), with the crossing pushed out to between $780$ and $1\,300$ on
uniform matrices (gradual) and to $811$ on integer ones; 16 is $1.2\times$
slower than naive at $n=32$; 32 and 64 are equivalent within noise at the top
of the sweep; 128 and 256 forgo the advantage below 512 and are $30$--$35\%$
slower than cutoff 32 at the top. Integer-valued matrices behave the same as uniform real ones to within
noise, as they should for an algorithm whose work does not depend on the
values.

The second baseline is the cautionary example. The textbook $i$-$j$-$k$ loop
order, which walks the second operand down a column with stride $n$, has a
fitted exponent of $4.3$ over $n\in[229, 2048]$---it is not cubic in practice,
for reasons we take up in Section~\ref{sec:mechanisms}---and at $n=2048$ it
takes $84$\,s where $i$-$k$-$j$ takes $5.0$\,s and Strassen $1.9$\,s. Swapping
two loop headers in the \emph{simple} algorithm is worth $17\times$ at that
size; Strassen's exponent, applied on top, is worth another $2.6\times$.

\subsection{Integer multiplication: Karatsuba vs.\ schoolbook}

\begin{figure}[t]
  \centering
  \includegraphics[width=0.85\linewidth]{fig_heatmap_bigint}
  \caption{$\log_2$ of the runtime ratio Karatsuba/schoolbook over operand
  length (32-bit limbs) and cutoff.}
  \label{fig:heatmap-bigint}
\end{figure}

Karatsuba is the cleanest case. With a cutoff of 32 limbs it overtakes
schoolbook multiplication at 34 limbs (interval $[32, 42]$; roughly $1\,100$
bits), essentially as soon as it does anything at all, and the measured
exponents over the top decade are $1.55$--$1.60$ for Karatsuba against
$1.97$--$2.01$ for schoolbook, within $0.04$ of $\log_2 3$ and $2$. Schoolbook
exceeded the 90-second budget at $4.5\times10^5$ limbs; at the last common
size, $3\times10^5$ limbs ($9.7$\,Mbit operands), Karatsuba is $18\times$
faster ($3.6$\,s vs $65$\,s). The crossing at 34 32-bit limbs is 17 64-bit
limbs (about 1\,090 bits). GMP's tuned parameter file for Alder Lake, the
previous generation of the same core design, sets
\texttt{MUL\_TOOM22\_THRESHOLD} to 13 64-bit limbs (830
bits)~\cite{gmp2022alderlake}, so the two land within a factor of $1.3$ of
each other. The agreement should not be over-read: GMP measures its threshold
against a hand-written assembly basecase and with an assembly Karatsuba, so
both sides of its comparison differ from ours, and the near-coincidence says
mainly that the crossover is insensitive to the quality of the
implementation when both sides are implemented with equal care.

With all-ones limbs---maximal carry propagation---schoolbook is unchanged but
Karatsuba is $1.8\times$ faster than on random limbs, because the sub-products
of the sum terms have predictable carries; the crossing is essentially unchanged (Karatsuba is ahead from the
first size at which it differs from schoolbook). The
cutoff sensitivity is asymmetric (Figure~\ref{fig:heatmap-bigint}): cutoffs of
32, 64 and 128 are indistinguishable at large $n$, 16 makes the crossing
gradual (the two are within noise from 34 to about 60 limbs), and 8 pushes it out
to $\approx700$ limbs and costs a factor of $1.9$ at every size above
($42.8$\,s vs $22.7$\,s at $10^6$ limbs). A cutoff four times too large costs
nothing; four times too small costs $2\times$.

\subsection{Shortest paths: Fibonacci heap vs.\ binary heap}

\begin{figure}[t]
  \centering
  \includegraphics[width=0.72\linewidth]{fig_ratio_dijkstra_by_dist}
  \caption{Runtime ratio Fibonacci-heap/binary-heap (decrease-key) Dijkstra
  by graph family (random graphs: average out-degree 16). Only the complete
  graph reaches the line, and only at the largest size that fits in memory.}
  \label{fig:ratio-dijkstra}
\end{figure}

The Fibonacci heap is the pair that does not cross. On random graphs of
average degree 2, 4, 16, 64 and 256 it trails the indexed binary heap by
$2.1\times$, $1.7\times$, $1.4\times$, $1.3\times$ and $1.3\times$ at the
largest sizes measured ($8.4\times10^6$ vertices at degrees 2--16, $134$\,M
edges at degree 16; $4.2\times10^6$ at degree 64 and $1.0\times10^6$ at 256), on 2-D grids by $2.6\times$ (in line with the $3.15\times$ that
Larkin, Sen and Tarjan measured against an implicit 4-ary heap on the USA road
network~\cite{larkin2014backtobasics}), and the gap narrows only slowly with
$n$ (Figure~\ref{fig:ratio-dijkstra}). The reason is in the operation counts
the harness logs: Dijkstra performs $0.05$ decrease-keys per vertex on
degree-2 graphs, $0.3$ on degree-4 graphs and grids, and $1.4$ at degree 16,
so the $\bigO(1)$ amortised decrease-key that justifies the Fibonacci heap is
exercised far less often than the $\bigO(\log n)$ delete-min that both heaps
share; Section~\ref{sec:mechanisms} argues that the Fibonacci heap's
delete-min is the more expensive of the two.

On complete graphs ($m = n(n-1)$, $6$ decrease-keys per vertex) the ratio
falls steadily from $2.4$ at 100 vertices to $1.0$ at $6\,500$ and $0.98$ at
$7\,700$ vertices ($60$\,M edges, the largest that fits in memory). This is
the regime $m \gg n\log n$ in which $\bigO(m+n\log n)$ should beat
$\bigO(m\log n)$, and the trend is as theory predicts---but what it reaches
within the measurable range is parity, not a win, and the atlas reports it as
such.

Even that parity is against the wrong baseline. The lazy variant,
\texttt{std::priority\_queue} with stale entries skipped on pop, does no
decrease-keys at all and is the fastest of the three on every random graph and
on the complete graph: $78$\,ms against $135$\,ms (Fibonacci) and $138$\,ms
(indexed binary) at $7\,700$ dense vertices; $6.3$\,s against $15.0$\,s and
$8.6$\,s at $8.4$\,M vertices of degree 4---the same ``at least $1.4\times$'' advantage for
dropping decrease-key that Chen et al.\ report on random graphs of average
degree 8~\cite{chen2007priority}. Its heap holds duplicate entries and does
more pops, but each pop touches a contiguous array and the extra work costs
less than any form of decrease-key bookkeeping. Only on grids, where the
frontier is small and duplicates accumulate, does the indexed binary heap edge
it out ($1.89$\,s vs $2.28$\,s). The Fibonacci heap never beats the lazy heap
in any cell, trailing it by $1.7\times$ even on the complete graph.

\subsection{Summary}

\begin{table}[t]
  \centering\scriptsize
  \caption{Sustained crossover points with 95\% bootstrap intervals (shortest
  interval over 1000 resamples of the repetitions). Sharp crossings are given
  as a point; gradual ones (interval wider than a grid step or a factor 1.5)
  as the range over which the two algorithms are within noise of each other;
  ``parity only'' marks a crossing that is never more than 10\% ahead by the
  largest size measured. $\rho(n_{\max})$ is the runtime ratio at the largest
  size measured. For matrices the crossing is reported on the
  power-of-two-rich grid and on all sizes separately. This table shows the
  default parameter of each pair (11-bit digits, cutoff 32); every
  parameter is in Appendix~\ref{app:table}.}
  \label{tab:crossovers}
  \input{../figures/table_crossovers}
\end{table}

Three of the four theoretically faster algorithms win, and early: at
$1\,500$ keys, 32 rows, and 34 limbs---and for the well-tuned configurations
these are sharp crossings with intervals of $\pm10\%$. Each has conditions:
radix sort needs the input not to be nearly or reverse sorted, Strassen needs the matrix
dimension to be rich in factors of two or the implementation to pad only
once, Karatsuba needs its cutoff not to be too small. The fourth reaches
parity on complete graphs at the edge of memory and loses everywhere else,
including to the lazy idiom that practitioners actually write.


\section{Mechanisms}\label{sec:mechanisms}

The reported run was made under WSL2, which hides the hardware performance
counters, so what follows is not measured: it is the set of explanations we
find most consistent with the timings, the working-set sizes and the
operation counts the harness logs. Each is stated as an attribution and could
be tested directly with the counter support the harness already has on
bare-metal Linux (\texttt{figures/fig\_counters\_*} are produced automatically
from such a run). 

\paragraph{Sorting: passes over memory against passes with unpredictable
branches.} LSD radix sort with $d$-bit digits makes $\lceil k/d\rceil$ passes
over the array for $k$-bit keys, each a sequential read and a scatter of
writes into $2^d$ output streams, with no data-dependent branches.
Introsort makes $\Theta(\log n)$ partition passes; each reads sequentially
from both ends of its range, which is cache-friendly, but the comparison at
every element is a branch that mispredicts about half the time on random
keys, and the number of passes grows with $n$. We attribute radix sort's
widening advantage above the crossing (Figure~\ref{fig:ratio-sorting}) to
this difference in pass count and branch behaviour, and the flat floor below
$\sim10^3$ keys to the fixed cost of the $2^d$-entry histograms. The
structured families support the branch explanation directly: on
nearly-sorted and reverse-sorted input introsort's comparisons become
predictable and it runs $6$--$11\times$ faster than on random keys of the
same size, while radix sort, whose work does not depend on the key order,
does not change. Whether radix sort's scattered stores are a bottleneck at
large $n$ is exactly what a counter run would show; the 16-bit-digit
results, where a $512$\,KB histogram exceeds L1d and the advantage shrinks,
suggest that the histogram's locality matters more than the output scatter.

\paragraph{Matrix multiplication: the access pattern of the base case.}
Strassen and Karatsuba change the number of multiplications but not how the
leaves access memory; their overhead is a constant number of extra additions
and temporaries per level. This is consistent with crossings that are small
and, on well-shaped inputs, insensitive to $n$ above the cutoff, and with
measured exponents within $0.04$ of theory for both recursive algorithms at
every cutoff. The $i$-$j$-$k$ baseline is the clearest case of an access
pattern, not an operation count, setting the time: its inner loop walks the
second operand down a column with stride $8n$ bytes, so computing one row of
the result touches every cache line of that operand; the operand is $8n^2$
bytes, which passes L1d (48\,KB) at $n\approx77$ and L2 (2\,MB) at $n=512$. The fitted exponent of $4.3$ over $n\in[229,2048]$ is consistent
with the operand leaving L2 within that range and every inner-loop step then
paying a cache miss. We attribute the odd-size penalty of Strassen to the
per-level padding copies, which the implementation makes explicitly and
whose count depends on the number of odd dimensions met in the recursion.

\paragraph{Shortest paths: pointer chasing against arrays.} A binary heap of
$n$ entries is one contiguous array and a sift costs $\log_2 n$ steps whose
addresses are computable in advance. A Fibonacci heap's consolidation walks a
linked root list of $\bigO(\log n)$ trees whose nodes were allocated in
insertion order, each step a dependent load. Both do $n$ delete-mins; the
Fibonacci heap pays this cost on every one and, on sparse graphs, recoups
little through decrease-keys because there are few: $0.05$ per vertex at
degree 2, $0.3$ at degree 4 and on grids, $1.4$ at degree 16 (operation
counts from the harness log). On complete graphs there are $6$ per vertex,
and the ratio falls towards 1 as theory predicts. This is the attribution
Larkin, Sen and Tarjan~\cite{larkin2014backtobasics} made from measured L1
miss rates on a road network, and our timings are consistent with it; we
have not measured it here.

\section{Discussion}\label{sec:discussion}

\paragraph{Crossovers are early, when they happen at all.}
For three of the four pairs the asymptotically faster algorithm wins from
inputs that fit comfortably in L1 cache: 1\,500 keys, a $32\times32$ matrix,
a 1\,100-bit integer. The common belief that ``the fancy algorithm only
pays off for huge inputs'' is, for these three, off by several orders of
magnitude. What the belief gets right is the fourth pair: an algorithm whose
advantage lies in an operation that realistic inputs rarely perform, and whose
implementation trades array traversal for pointer chasing, can fail to win at
any size that fits in memory.

\paragraph{The parameter moves the crossover more than the exponent does.}
For Karatsuba the base-case cutoff shifts $n^*$ by a factor of twenty (34
limbs at cutoff 32, 700 at cutoff 8) and halves the asymptotic advantage; for
Strassen it moves $n^*$ between 32 and roughly 1\,000 and changes the speed at
$n=4096$ by $1.6\times$; the shape of the matrix---how many factors of two
its dimension has---moves Strassen's crossing by a factor of ten; and the loop
order of the \emph{baseline} is worth $17\times$ at $n=2048$ where
Strassen's exponent is worth $2.6\times$. In the measured range, engineering decisions that the
asymptotic analysis cannot see dominate the decision it can. The practical
reading is that a recursive algorithm is only as good as its base-case policy,
and that a ``naive'' baseline must be a competent one before any comparison
means anything.

\paragraph{Wide intervals are answers, not failures.}
Several crossings in Table~\ref{tab:crossovers} have intervals spanning a
factor of two or more: Strassen with a too-small cutoff, Karatsuba with a
too-small cutoff, radix sort with too-wide digits. In each the ratio curve approaches 1 slowly, and the
bootstrap faithfully reports that the crossing could be anywhere in a wide
range. For a practitioner this is the useful fact: over that range the two
implementations are within noise of each other and the simpler one should be
preferred for its simplicity; the asymptotic winner only earns its complexity
beyond the top of the interval, and then only by the modest ratio we report
there. Sharp crossings, by contrast, are where the decision is worth making
carefully, and the sorting and Karatsuba (cutoff 32) rows show the grid can
resolve them to $\pm10\%$ and to a few limbs respectively.

\paragraph{Input structure can flip the verdict.}
Radix sort's loss on nearly-sorted and reverse-sorted input and the Fibonacci
heap's parity only on complete graphs are the same phenomenon seen from two
sides. An algorithm
with a worse bound but adaptive behaviour (introsort, the lazy heap) exploits
structure that the input happens to have; an algorithm with a better bound
and oblivious behaviour (radix sort, the Fibonacci heap's decrease-key) does
exactly the same work regardless. Whether that is a feature or a bug depends
on whether the structure is present, and the atlas gives the practitioner a
way to check: the key is the input family, not the size.

\paragraph{Measured exponents are honest in the middle and lie at the edges.}
Over the largest decade of each sweep the fitted exponents sit close to theory
for the compute-bound algorithms (Karatsuba 1.55--1.60, schoolbook
1.97--2.01, Strassen 2.77--2.86) and drift above it for the memory-bound ones (introsort 1.14 on 64-bit keys where $n\log n$ predicts $\approx1.06$
over that decade; the $i$-$j$-$k$ loop 4.3 where theory says 3). A fitted exponent on a log--log plot is a statement about the cache
hierarchy as much as about the algorithm, which is why the cache boundaries
are drawn on every runtime figure.

\paragraph{What the atlas cannot say.}
These are single-threaded scalar implementations on one machine, measured
under a hypervisor without hardware counters. A vectorised radix sort or a
BLAS kernel would move every crossover, and not necessarily in the direction one
expects, since vectorisation may help the simple kernel more than the
recursive one built on it. The harness is
offered as the way to redraw the map rather than the map as a permanent
answer; Section~\ref{sec:repro} gives the one command that does so.

\paragraph{Guidance.}
For a practitioner choosing between these pairs today, on hardware like
ours:
\begin{itemize}\itemsep0pt
  \item Sort with radix sort (11-bit digits) when keys are integers, inputs
  exceed about 1\,500 elements, and the data is not already nearly or
  reverse sorted; use \texttt{std::sort} otherwise.
  \item Use Strassen with a base case of 32--64 for dense matrices above
  about 64 rows when the dimension is a multiple of a high power of two,
  or above about 300 otherwise (or pad once at the top)---and fix the loop
  order of the baseline first.
  \item Use Karatsuba above about 32 limbs with a cutoff of 32--128; the
  libraries already do.
  \item Do not use a Fibonacci heap for Dijkstra on any graph that fits in
  memory. Use a binary heap with lazy deletion unless the graph is a grid-like
  network, where an indexed binary heap is slightly better.
\end{itemize}


\section{Threats to validity}
\paragraph{Single machine.} Crossovers depend on cache sizes, memory bandwidth
and the branch predictor. We report one machine in full and treat
cross-machine generalisation as future work, offering the harness as the means.
\paragraph{Implementation quality.} Both sides are deliberately plain. A SIMD
radix sort or a BLAS kernel would move every crossover, and so, to a smaller
degree, would compiling with \texttt{-O3 -march=native}; the harness provides
that variant (\texttt{make native-quick}) but we have not run it for this
paper.
\paragraph{Virtualisation.} WSL2 adds timer and scheduling noise and hides
counters. The interleaving and the bootstrap are designed to make the ratio
robust to it, and the interval widths make the residual noise visible rather
than hidden.
\paragraph{Input families.} Each family is a convenience sample of structures
known to matter (sortedness, key width, density); the atlas is a map, not a
census.


\section{Reproducibility}\label{sec:repro}
All code, raw measurements, analysis scripts and figure sources are at
\url{https://github.com/dhanushp08/algorithm-crossover-atlas} (tag \texttt{v1.0}). \texttt{make} builds and tests the code;
\texttt{make full} and \texttt{make refine} re-measure the sweeps, and the
analysis command in the README regenerates every number, table and figure in
this paper from the raw measurements; \texttt{docs/BENCHMARKING.md} describes
the machine setup. The raw data of the reported runs (\texttt{full},
\texttt{refine} and \texttt{patch}, Section~\ref{sec:method-patch}, together
with three preliminary \texttt{quick} runs) is in \texttt{results/}.

\section{Conclusion}
Asymptotic analysis answers ``which algorithm wins eventually''; this paper
answers ``from what size'' for four pairs that programmers choose between every
day, on one well-described machine, with intervals rather than point guesses.
Three of the four crossovers are small enough to fit in first-level cache
and, for well-tuned configurations, sharp; one does not exist on any sparse
input and reaches only parity on dense ones; and in every case the tuning
parameter, the input structure or the shape of the problem moved the answer
more than the exponent did. Where the data cannot pin a crossover down, we report the range
rather than a point, because within that range the choice is immaterial.
The harness that produced the atlas rebuilds it with one command, so the map
can be redrawn as the hardware underneath it changes.

\bibliographystyle{plainnat}
\bibliography{references}

\clearpage
\appendix
\section{All crossover points}\label{app:table}
{\scriptsize
\input{../figures/table_crossovers_full}
}

\end{document}
